\seccionInforme{1. Descripción general del programa}
\phantomsection\addcontentsline{toc}{section}{1. Descripción general del programa}

\vspace{0.2cm}
\phantomsection\addcontentsline{toc}{subsection}{1.1 Contexto y propósito}
\colabtitle{1.1 Contexto y propósito}

\textbf{MatriLab EPIS} es una aplicación móvil para gestionar la matrícula en los grupos de laboratorio de la Escuela Profesional de Ingeniería de Sistemas de la UNSA. Es el mismo dominio que el equipo trabajó en el LAB 01 (selección de herramienta y requisitos) y en el LAB 02 (prototipo RMAD en FlutterFlow). En esta práctica la aplicación se reconstruye con un \textbf{framework multiplataforma}, \textbf{React Native}, usando únicamente los componentes propios de la plataforma para organizar cada interfaz.

El programa atiende a tres usuarios. El \textbf{estudiante} consulta la oferta de grupos con un semáforo de vacantes, se matricula sin cruces de horario y obtiene un pase digital con código QR. El \textbf{docente o jefe de prácticas} revisa la nómina de sus grupos y registra la asistencia leyendo el QR con la cámara del teléfono. La \textbf{coordinación de laboratorios} define la capacidad de cada ambiente y la ventana de matrícula. Las doce pantallas comparten un mismo código que se ejecuta como aplicación nativa en Android y como aplicación web.

Los datos de esta versión son de demostración y viven en memoria: el objetivo de la práctica es el diseño y la organización de las interfaces. Las reglas de negocio (cruce de horarios, aforo, ventana de matrícula y verificación del pase) sí son reales y se verifican con pruebas automatizadas.

\vspace{0.2cm}
\phantomsection\addcontentsline{toc}{subsection}{1.2 Árbol de problemas}
\colabtitle{1.2 Árbol de problemas}

El diseño parte del diagnóstico del LAB 01 sobre el proceso de matrícula de laboratorios. Se aplicó la técnica del árbol de problemas de la metodología del marco lógico [17]: un problema central, sus causas directas e indirectas y los efectos que produce (Figura \ref{fig:arbol_problemas}).

\begin{figure}[!htbp]
\centering
\safeimage{1.0\textwidth}{img/diagramas/arbol_problemas.png}{Árbol de problemas del proceso de matrícula de laboratorios de la EPIS-UNSA}
\caption{Árbol de problemas del proceso de matrícula de laboratorios de la EPIS-UNSA}
\label{fig:arbol_problemas}
\end{figure}

El problema central es que \textbf{la matrícula en los laboratorios es lenta, poco transparente y genera sobrecupos y cruces de horario}. Las causas de fondo son listas que no se actualizan en tiempo real, la falta de contraste entre el grupo práctico y la teoría inscrita, inscripciones simultáneas sin una operación atómica, la ausencia de una constancia digital y la coordinación de las salas por correo. Los efectos llegan al aula: rechazos al final del trámite, permutas cuando las clases ya empezaron, salas con más estudiantes que computadoras y nóminas tardías.

\vspace{0.2cm}
\phantomsection\addcontentsline{toc}{subsection}{1.3 Árbol de objetivos y alcance}
\colabtitle{1.3 Árbol de objetivos y alcance}

Cada condición negativa se convirtió en la situación deseada (Figura \ref{fig:arbol_objetivos}). Los \textbf{medios fundamentales}, en la base del árbol, son acciones concretas que el software puede ejecutar; por eso se tomaron como origen de las funcionalidades. Debajo de cada medio se indican las interfaces de MatriLab que lo implementan.

\begin{figure}[!htbp]
\centering
\safeimage{1.0\textwidth}{img/diagramas/arbol_objetivos.png}{Árbol de objetivos: los medios fundamentales se traducen en interfaces de MatriLab}
\caption{Árbol de objetivos: los medios fundamentales se traducen en interfaces de MatriLab}
\label{fig:arbol_objetivos}
\end{figure}

El alcance de la práctica cubre los cinco medios con \textbf{doce interfaces} (I1 a I12). La autenticación real con Google Workspace y la persistencia en la nube, que en el LAB 02 se resolvieron con Firebase, se simulan aquí con un inicio de sesión de demostración y un estado local, de modo que cualquier evaluador puede probar la aplicación sin cuentas ni conexión.

\vspace{0.2cm}
\phantomsection\addcontentsline{toc}{subsection}{1.4 Actores y requisitos atendidos}
\colabtitle{1.4 Actores y requisitos atendidos}

\begin{table}[H]
\centering\small
\caption{Actores del sistema y lo que cada uno puede hacer en la aplicación}\label{tab:actores}
\renewcommand{\arraystretch}{1.25}
\begin{tabular}{|>{\raggedright\arraybackslash}p{\dimexpr 0.170\linewidth-2\tabcolsep-2\arrayrulewidth\relax}|>{\raggedright\arraybackslash}p{\dimexpr 0.220\linewidth-2\tabcolsep-2\arrayrulewidth\relax}|>{\raggedright\arraybackslash}p{\dimexpr 0.460\linewidth-2\tabcolsep-2\arrayrulewidth\relax}|>{\raggedright\arraybackslash}p{\dimexpr 0.150\linewidth-2\tabcolsep-2\arrayrulewidth\relax}|}
\hline
\rowcolor{colabpyblue} \color{white}\textbf{Rol} & \color{white}\textbf{Cuenta de demostración} & \color{white}\textbf{Responsabilidades en MatriLab} & \color{white}\textbf{Interfaces} \\ \hline
Estudiante & Lopez Arela, Ower Frank · CUI 20213412 & Consultar la oferta, matricularse sin cruces, ver su horario, mostrar su pase QR y anular una matrícula dentro de la ventana. & I1–I6, I12 \\ \hline
\rowcolor{colabpybg} Docente / jefe de prácticas & Mg. Carla Medina Rojas (ficticia) & Revisar la ocupación de sus grupos, consultar la nómina, registrar asistencia con QR o código manual y exportar la nómina. & I1, I7–I9, I12 \\ \hline
Coordinación de laboratorios & \texttt{laboratorios.\allowbreak{}epis@unsa.\allowbreak{}edu.\allowbreak{}pe} & Ajustar capacidad y estado de cada ambiente, definir la ventana de matrícula y revisar los grupos agotados. & I1, I10–I12 \\ \hline
\end{tabular}
\end{table}

\begin{table}[H]
\centering\small
\caption{Requisitos funcionales del LAB 01 y cómo se resuelven en la aplicación}\label{tab:rf}
\renewcommand{\arraystretch}{1.25}
\begin{tabular}{|>{\raggedright\arraybackslash}p{\dimexpr 0.090\linewidth-2\tabcolsep-2\arrayrulewidth\relax}|>{\raggedright\arraybackslash}p{\dimexpr 0.250\linewidth-2\tabcolsep-2\arrayrulewidth\relax}|>{\raggedright\arraybackslash}p{\dimexpr 0.510\linewidth-2\tabcolsep-2\arrayrulewidth\relax}|>{\raggedright\arraybackslash}p{\dimexpr 0.150\linewidth-2\tabcolsep-2\arrayrulewidth\relax}|}
\hline
\rowcolor{colabpyblue} \color{white}\textbf{Código} & \color{white}\textbf{Requisito} & \color{white}\textbf{Cómo se resuelve} & \color{white}\textbf{Interfaz} \\ \hline
RF-01 & Autenticación institucional & Validación del dominio @unsa.edu.pe y de la contraseña; acceso de demostración por rol. & I1 \\ \hline
\rowcolor{colabpybg} RF-02 & Catálogo de cursos y grupos & Lista agrupada por curso con buscador y filtros. & I3 \\ \hline
RF-03 & Vacantes en vivo & Semáforo y barra de ocupación que se recalculan al matricular o anular. & I2, I3, I4 \\ \hline
\rowcolor{colabpybg} RF-04 & Validación anti-cruces & Contraste del grupo con los laboratorios y la teoría inscritos (RN-02). & I4 \\ \hline
RF-05 & Reserva atómica & El cupo se descuenta y la matrícula se registra en una sola operación. & I4 \\ \hline
\rowcolor{colabpybg} RF-06 & Ficha digital con QR & Pase con código QR firmado y código de verificación de 8 caracteres. & I5 \\ \hline
RF-07 & Horario consolidado y nómina & Cuadrícula semanal con teoría y laboratorios; nómina por grupo. & I6, I7, I8 \\ \hline
\rowcolor{colabpybg} RF-08 & Asistencia por escaneo QR & Cámara del teléfono con lector QR o ingreso manual del código. & I9 \\ \hline
RF-09 & Parametrización de laboratorios & Capacidad, estado y datos de cada sala; ventana de matrícula validada. & I10, I11 \\ \hline
\rowcolor{colabpybg} RF-10 & Exportación de nóminas & CSV con separador punto y coma compartido con la hoja nativa del sistema. & I8 \\ \hline
\end{tabular}
\end{table}

\vspace{0.2cm}
\phantomsection\addcontentsline{toc}{subsection}{1.5 Elección del framework}
\colabtitle{1.5 Elección del framework}

La guía permite Flutter, React Native o Xamarin. El equipo eligió \textbf{React Native 0.88 con Expo SDK 58} por tres razones. Primero, React Native dibuja la interfaz con \textbf{vistas nativas} de cada sistema operativo: un \texttt{View} es un \texttt{View\allowbreak{}Group} en Android y un \texttt{UIView} en iOS [2], lo que responde directamente al requisito de usar los componentes propios de la plataforma. Segundo, el LAB 02 ya se había construido sobre Flutter mediante FlutterFlow, y usar otro framework permite comparar ambos enfoques con evidencia propia. Tercero, Expo [14] permite ejecutar la aplicación en un emulador o teléfono con Expo Go sin compilar proyectos nativos y exportarla a la web con el mismo código. Se complementó con TypeScript, React Navigation 7 para la navegación nativa, \texttt{expo-camera} para el lector QR y Jest para las pruebas.

